\documentclass[10pt,a4paper]{amsart}
\usepackage[margin=19mm]{geometry}
\usepackage{amsmath,amssymb,mathtools}
\usepackage[scr=rsfs,cal=euler]{mathalfa}
\usepackage{xfrac}
\usepackage[unicode,hidelinks,bookmarks=false]{hyperref}
\newcommand{\ie}{i.e.\ }
\newcommand{\cf}{cf.\ }
\newcommand{\resp}{respectively\ }
\newcommand{\Subsection}{\subsection}
\providecommand{\nobreakdash}{\nobreak\hbox{-}}
\setlength{\emergencystretch}{2em}
\multlinegap0pt
\usepackage{amssymb}
\usepackage{xcolor}
\usepackage[shortlabels]{enumitem}
\newtheorem{prop}{Proposition}
\newcommand{\D}{\mathcal{D}}
\renewcommand{\O}{\Omega}
\newcommand{\R}{\mathbb{R}}
\renewcommand{\tt}[1]{\text{#1}}
\title{$L^p$ stability of Vortex Patches in Two Dimensional Domains}
\author{Zelin Dong}
\date{Source: arXiv:2608.13998v1 (CC BY 4.0)}
\begin{document}
\maketitle

\noindent\emph{Source and licence.} $L^p$ stability of Vortex Patches in Two Dimensional Domains. Version arXiv:2608.13998v1 (CC BY 4.0), distributed by arXiv under the Creative Commons Attribution 4.0 International licence, http://creativecommons.org/licenses/by/4.0/, as recorded in the arXiv licence metadata for this exact version and shown on the abstract page https://arxiv.org/abs/2608.13998v1. Full licence notice: \texttt{licenses/arxiv-2608.13998-CC-BY-4.0.txt}.\par\smallskip

\noindent\textit{Editorial scope.} Counted unit: the article's prop prop:weak\_finite\_volume\_energy\_convergence (source0.tex lines 1200--1208). The article's complete original proof of it is delivered with it (source0.tex lines 1209--1283, one \texttt{proof} environment of 75 lines). No mathematics is written, repaired or re-derived: the statement and the proof are the article's own lines, in the article's own notation, and no retained line is substituted or re-flowed.  Retained uncounted as the source-local context the proof needs: lines 152--154; lines 286--292.  Nothing was omitted from any retained span, so the package carries no omission record.  No retained line contains a citation, so the wrapper carries no bibliography and no bibliography entry.  No retained line contains a figure environment, an image-inclusion command or drawing code, so no image is delivered and the package declares no asset dependency.  Editorial formatting only, in addition: an AMS \texttt{amsart} preamble that restates the article's own declarations and macros used by the retained lines, namely the environments \texttt{prop} on separate counters, together with the standard packages those lines need.  The retained environments print in this document as Proposition 1 in order of appearance, whereas the article prints the same items with the numbering of its own full document.  The complete native source of the article as distributed in the arXiv e-print for this exact version is bundled unchanged as \texttt{sources/207627/source0.tex}.
                \begin{align}
                    \int_{D_l} x_2^q dx<\infty, \tt{ for some } q\geq 0,\tt{ and } ~ \tt{Vol}(D_u)<\infty. \label{def:weak_finite_volume_domain_assumption}
                \end{align}

    \begin{align}
        &\int_{\Omega}G(x,y)\omega(y)dy\leq Cx_2^{\frac{p-1}{2p-1}}\|x_2 \omega\|_1^{\frac{p-1}{2p-1}}\|\omega\|_p^{\frac{p}{2p-1}},\label{ineq:estimate of stream fct}\\
        &\int_{\Omega}G(x,y)\omega(y)dy\leq C\|x_2 \omega\|_1^{\frac{2p-2}{3p-2}}\|\omega\|_p^{\frac{p}{3p-2}},\label{ineq:sup estimate of stream fct}\\
        &\int_{\Omega}\int_{\Omega}G(x,y)\omega_1(y)\omega_2(x)dxdy\leq C\|x_2\omega_1\|_1^{\frac{2p-2}{3p-2}}\|\omega_1\|_p^{\frac{p}{3p-2}}\|x_2\omega_2\|_1^{\frac{2p-2}{3p-2}}\|\omega_2\|_p^{\frac{p}{3p-2}},\label{ineq:estimate of energy crossing term}\\
        &E[\omega]\leq C\|x_2\omega\|_1^{\frac{4p-4}{3p-2}}\|\omega\|_p^{\frac{2p}{3p-2}},\label{ineq:estimate of kinetic energy}\\
        &|E[\omega_1]-E[\omega_2]|\leq C\|x_2(\omega_1-\omega_2)\|_1^{\frac{2p-2}{3p-2}}\|\omega_1-\omega_2\|_p^{\frac{p}{3p-2}}\|x_2(\omega_1+\omega_2)\|_1^{\frac{2p-2}{3p-2}}\|\omega_1+\omega_2\|_p^{\frac{p}{3p-2}}.\label{ineq:estimate of energy difference}
    \end{align}

\begin{prop}[Convergence of the kinetic energy]
\label{prop:weak_finite_volume_energy_convergence}
Let $\O=\D$ and one of the following holds: (1), $q\in[0,2]$, $p>1$; (2), $q\in (2, + \infty)$, $ 1 < p \leq q/(q-1)$.
Suppose \(\{\omega_n\} \subset L^p(\D)\) and \(\omega \in L^p(\D)\), such that 
\begin{align*}
    \omega_n \rightharpoonup \omega  \tt{ in } L^p(\D).
\end{align*}
If each $\omega_n$ is nonnegative, then \(E[\omega_n] \to E[\omega]\).
\end{prop}

\begin{proof}
    \underline{\textit{Case (1):}} It is sufficient to prove that $G\in L^{p^\prime}(\D \times \D)$ for any $1/p+1/p^{\prime}=1$. This follows form:
    \begin{align*}
        & \int_{\D}\int_{\D}G^{p^{\prime}}(x,y)dxdy 
        % \leq \int_{\D}\int_{\D}G_H^{p^{\prime}}(x,y)dxdy, 
        \\
        \leq & \int_{D_u}\int_{D_u}G_H^{p^{\prime}}(x,y)dxdy + 2\int_{D_l}\int_{D_u}G_H^{p^{\prime}}(x,y)dxdy + \int_{D_l}\int_{D_l}G_H^{p^{\prime}}(x,y)dxdy, \\
        \leq & \int_{D_u}\int_{D_u}G_H^{p^{\prime}}(x,y)dxdy + 3\int_{D_l}\int_{\R^2_+}G_H^{p^{\prime}}(x,y)dxdy, \\
        \leq & \int_{D_u}\int_{D_u}G_H^{p^{\prime}}(x,y)dxdy + C\int_{D_l}y_2^2dy.
    \end{align*}
    The first term is finite because $\tt{Vol}(D_u)<\infty$, and the last term is finite by $q\in[0,2]$ in the sense that 
    \begin{align*}
        \int_{D_l}y_2^2 dy \leq \int_{D_l}y_2^q dy <\infty.
    \end{align*}
    
    \underline{\textit{Case (2):}} We first demonstrate the concentration of $\{x_2\omega_n\}$ in the sense that
    \begin{align}
        \sup_{n}\int_{D_l \cap \{|x_1|>N\}} x_2 \omega_n dx\leq \epsilon. \label{ineq:concentration of x_2 omega_n}
    \end{align}
    Suppose \eqref{ineq:concentration of x_2 omega_n} fails. Then there exists $\epsilon_0>0$ and a subsequence $\{\omega_{n_k}\}$ such that for every $k \in \mathbb{Z}_{\geq 1}$,
    \begin{align*}
        \int_{D_l \cap \{|x_1|>k\}} x_2 \omega_{n_k} dx > \epsilon_0.
    \end{align*}
    Thus 
    \begin{align*}
        \epsilon_0 < \int_{D_l \cap \{|x_1|>k\}} x_2 \omega_{n_k} dx \leq  \left(\int_{D_l \cap \{|x_1|>k\}} x_2^{p^\prime} dx\right)^{1/p^\prime}\left(\int_{D_l \cap \{|x_1|>k\}} \omega_{n_k}^p dx\right)^{1/p}.
    \end{align*}
    where $p^\prime=p/(p-1) $. The assumption $p \leq q/(q-1)$ implies $p^\prime\geq q$, and the weak finite volume condition \eqref{def:weak_finite_volume_domain_assumption} on $D_l$ gives
    \begin{align*}
        \int_{D_l}x_2^{p^\prime}dx\leq \int_{D_l}x_2^{q}dx <+\infty.
    \end{align*}
    Then 
    \begin{align*}
        \int_{D_l \cap \{|x_1|>k\}} x_2^{p^\prime} dx \rightarrow 0, \tt{ as } k \rightarrow \infty,
    \end{align*}
    we conclude
    \begin{align*}
        \|\omega_{n_k}\|_{p} \geq (\int_{D_l \cap \{|x_1|>k\}} \omega_{n_k}^p dx)^{1/p} \geq \epsilon_0 \Big(\int_{D_l \cap \{|x_1|>k\}} x_2^{p^\prime} dx\Big)^{-1/p^\prime}\rightarrow +\infty.
    \end{align*}
    $\{\omega_{n_k}\}$ is uniformly bounded in $L^p$ because $\omega_{n_k}\rightharpoonup\omega$ in $L^p(D)$, a contradiction arises. 
    To complete the proof, fix $0<\epsilon<1$ and define
    $$
        D_0=D_u\cup(D_l \cap \{|x_1|\leq N\}),
    $$
    where $N=N(\epsilon^{\frac{3p-2}{2p-2}})$ is the constant obtained in the previous step. By properly increasing $N$, if needed, we can assume
    \begin{align*}
        \int_{\D\backslash D_0} x_2\omega dx\leq \epsilon^{\frac{3p-2}{2p-2}}.
    \end{align*}
    Then 
    \begin{align*}
        |E[\omega_n] & - E[\omega]| \leq |E[\omega_n|_{D_0}]-E[\omega|_{D_0}]| \\
        & + \int_{\mathcal{D}}\int_{\mathcal{D}} G(x,y)\omega_n|_{D_0}(x)\omega_n|_{\mathcal{D}\backslash D_0}(y)dydx + E[\omega_n|_{\mathcal{D}\backslash D_0}]\\
        & + \int_{\mathcal{D}}\int_{\mathcal{D}} G(x,y)\omega|_{D_0}(x)\omega|_{\mathcal{D}\backslash D_0}(y)dydx + E[\omega|_{\mathcal{D}\backslash D_0}].
    \end{align*}

    For the last four terms, using \eqref{ineq:estimate of kinetic energy}, \eqref{ineq:estimate of energy crossing term} together with $$\sup_n\|x_2\omega_n|_{\mathcal{D}\backslash D_0}\|_1, \|x_2 \omega|_{\mathcal{D}\backslash D_0}\|_1\leq \epsilon^{\frac{3p-2}{2p-2}},$$ 
    we obtain
    \begin{align*}
        |E[\omega_n] & - E[\omega]| \leq |E[\omega_n|_{D_0}]-E[\omega|_{D_0}]| + C M^2\epsilon.
    \end{align*}
    where $M$ is a large constant such that $\sup_n\|\omega_n\|_{p}, \|\omega\|_{p}\leq M$.

    For the first term, note that $\tt{Vol}(D_0)<\infty$, so $G(x,y) \in L^{p^{\prime}}(D_0 \times D_0)$ with $1/p+1/{p^\prime}=1$ and $\omega_n(x)\omega_n(y) \rightharpoonup \omega(x)\omega(y)$ in $L^p(D_0 \times D_0)$. Hence, as $n\to\infty$
    \begin{align*}
        E[\omega_n|_{D_0}] = & \int_{D_0}\int_{D_0} {G}(x,y)\omega_n(x)\omega_n(y)dydx \\
        & \longrightarrow \int_{D_0}\int_{D_0} {G}(x,y)\omega(x)\omega(y)dydx = E[\omega|_{D_0}].
    \end{align*}

    In conclusion, as $n$ is sufficiently large, 
    \begin{align*}
        |E[\omega_n] & - E[\omega]| \leq 2CM^2\epsilon.
    \end{align*}
    Since $\epsilon>0$ can be arbitrary small, the proof is complete.
    
\end{proof}

\end{document}
